\subsection{The fixed-lift compactness question in Task F.4}

\begin{statement}{F.20}{Proved}{The finite coboundary criterion in the profinite kernel}
Use the data of Task F.4: $K/\mathbb Q_p$ is finite,
$\Delta=G_{K(\mu_\infty)}$, $\pi$ is the full geometric
profinite fundamental group, and
$q:\pi\twoheadrightarrow\pi^{(p')}$ has kernel $N_p$.
All actions are the genuine continuous actions induced
by the given section $s$.
For a fixed $\bar h\in(\pi^{(p')})^{s(\Delta)}$,
choose $h\in\pi$ above it and put
$d_\sigma=h^{-1}s(\sigma)h s(\sigma)^{-1}$.
Then $\bar h$ has a fixed lift in $\pi$ exactly when
the image of $d$ is a coboundary in $N_p/U$ for every
member $U$ of a cofinal directed family of open normal
subgroups of $N_p$ stable under $s(\Delta)$.
Such a family exists.  The solution set is exactly
the inverse limit of the finite solution sets; no
additional derived-limit obstruction occurs.
\end{statement}
\begin{proof}
Write $\sigma(n)=s(\sigma)n s(\sigma)^{-1}$.
The fixed-image hypothesis puts $d_\sigma$ in $N_p$,
and multiplication gives
$d_{\sigma\tau}=d_\sigma\sigma(d_\tau)$.
Moreover $hn$ is fixed precisely when
\[
             d_\sigma=n\sigma(n)^{-1}
                         \quad\text{for every }\sigma.       \tag{F.20}
\]
These identities prove the cocycle assertion and its
exact relation to the desired lift, including the
side on which the correcting element is multiplied.

We justify the cofinal stable family even though
$N_p$ need not be finitely generated.
For an open normal subgroup $U\subset N_p$, joint
continuity of the action and compactness of $N_p$
give an open neighborhood of the identity in
$s(\Delta)$ acting trivially modulo $U$ on all of
$N_p$.  Concretely, for each point choose an action
neighborhood preserving its coset modulo $U$;
finitely many point neighborhoods cover $N_p$,
and intersect their action neighborhoods and
their inverses.  A smaller open subgroup then
stabilizes $U$.  The orbit of $U$ is finite, and
its intersection is an open normal subgroup
contained in $U$ stable under all of $s(\Delta)$.
Taking these intersections and finite common
refinements gives the claimed family.

For each such $U$ let
\[
 X_U=\{nU\in N_p/U:
      d_\sigma=n\sigma(n)^{-1}\bmod U\text{ for every }\sigma\}.
\]
This is a finite set and is exactly the finite-level
solution set.  Its inverse image $Y_U$ in the compact
group $N_p$ is closed.  Under the stated finite-level
condition it is nonempty.  A common refinement of
finitely many $U$ supplies a nonempty $Y$ contained
in all corresponding $Y_U$, so these closed sets
have the finite-intersection property.  Compactness
gives $n\in\bigcap_UY_U$.  Since $\bigcap_UU=1$,
equation (F.20) holds in $N_p$ for every $\sigma$.
Conversely an actual solution gives every finite
solution.  The identification
$N_p=\varprojlim_U N_p/U$ identifies the total
solution set with $\varprojlim_UX_U$.
This proves the exact criterion, not merely a
necessary collection of finite tests.
\end{proof}

\begin{statement}{F.22}{Proved}{A full free pro-prime Dehn-twist fixed group}
Let $F_\ell$ be the free pro-$\ell$ group on $x,y,z$ and let
$0\ne m\in\mathbb Z_\ell$.  The continuous automorphism
\[
 \alpha_m(x)=y^mxy^{-m},\qquad
 \alpha_m(y)=y,\qquad \alpha_m(z)=z
\]
has fixed subgroup exactly the closed free factor
$\langle y,z\rangle$.  This is a statement about the entire
free pro-$\ell$ group.
\end{statement}
\begin{proof}
We use the continuous Magnus map
\[
 F_\ell\longrightarrow
 1+(X,Y,Z)\mathbb F_\ell\langle\!\langle X,Y,Z\rangle\!\rangle,
 \qquad x\mapsto1+X,\ y\mapsto1+Y,\ z\mapsto1+Z.            \tag{F.22.1}
\]
Here the ring has its degree-adic topology, and the unit group
on the right is an inverse limit of finite $\ell$-groups.
The map is injective.  To check this rather than assume a
linearization, the inverse ring map sends $X,Y,Z$ to
$x-1,y-1,z-1$ in $\mathbb F_\ell[[F_\ell]]$.
These substitutions converge: the augmentation ideal in the
group algebra of every finite $\ell$-group is nilpotent.
The latter fact follows by induction using a central subgroup
of order $\ell$, whose generator $c$ satisfies $(c-1)^\ell=0$.
The two continuous ring homomorphisms are inverse, by their
values on the generators.  Finally $F_\ell$ embeds in the
unit group of its completed group algebra, because its image
in each finite group algebra retains the distinct group basis
elements.  This proves the assertion about (F.22.1).

Put $u(Y)=(1+Y)^m$ in $\mathbb F_\ell[[Y]]$.
It is nonconstant: writing $m=\ell^v m_0$ with $m_0$ a unit,
its first nonconstant term has degree $\ell^v$ and nonzero
coefficient.  The induced substitution on the ring is
$X\mapsto u(Y)Xu(Y)^{-1}$, with $Y,Z$ fixed.
For a word $w=w_1\cdots w_k$ in the two letters $X,Z$,
collect all coefficients of
\[
                  Y^{a_0}w_1Y^{a_1}\cdots w_kY^{a_k}
\]
as the single commutative series
$f_w(Y_0,\ldots,Y_k)$.  Every noncommutative series has this
unique coefficient decomposition.  The substitution multiplies
$f_w$ by
\[
        \prod_{i:\,w_i=X}u(Y_{i-1})u(Y_i)^{-1}.             \tag{F.22.2}
\]
For every $w$ containing $X$ this product differs from one:
at the left endpoint of the first consecutive block of $X$'s
the corresponding variable has exponent $+1$, with no inverse
factor to cancel it.  The commutative power-series ring over
$\mathbb F_\ell$ is an integral domain.  A fixed series therefore
has $f_w=0$ whenever $w$ contains $X$.

For a fixed group element, its Magnus series consequently uses
only $Y,Z$.  The continuous retraction
$F_\ell\to\langle y,z\rangle$ given by $x\mapsto1$ sets $X=0$
in its Magnus series and leaves that series unchanged.
Injectivity of (F.22.1) makes the element equal to its retraction.
Conversely $\alpha_m$ fixes the stated free factor pointwise.
\end{proof}

\begin{statement}{F.23}{Proved}{All prime-to-five pro-prime fixed quotients in the three local examples}
Let $\Delta=G_{\mathbb Q_5(\mu_\infty)}$ and let
$\pi^{(\ell)}$ denote the maximal geometric pro-$\ell$ quotient.
For each $\ell\ne5$, the first two examples of Task F.4 satisfy
\[
  \bigl(\pi^{(\ell)}\bigr)^{s(\Delta)}
       =\pi^{(\ell)}=F_\ell(2).
\]
For the third example, with the point $2$ on
$\mathbb P^1\setminus\{0,1,5,\infty\}$, choose peripheral
generators $a,b,c$ at $0,5,1$ based through the outer component
of its semistable model and put $d=ab$.  Then
\[
 \bigl(\pi^{(\ell)}\bigr)^{s(\Delta)}
                  =\langle d,c\rangle\simeq F_\ell(2).     \tag{F.23}
\]
The natural quotient $\pi\twoheadrightarrow\pi^{(\ell)}$
is continuous and $s(G_{\mathbb Q_5})$-equivariant.  It sends
the full fixed group into the displayed subgroup; surjectivity
of that restriction is not asserted.
\end{statement}
\begin{proof}
The smooth marked integral model in the first two examples,
together with its integral point or unit tangential base point,
gives based prime-to-five specialization.  Inertia acts trivially
on this quotient.  Since all finite unramified extensions of
$\mathbb Q_5$ are generated by prime-to-five roots of unity,
$\Delta$ is contained in inertia, proving these two assertions.

For the third example blow up the collision $x=5=0$ and use the
node chart $XV=5$, with $X=x$ and $V=5/x$.
The outer component has node $0$, cusps $1,\infty$, and the
smooth reduced base point $2$.  The bubble has node $0$ and
cusps $1,\infty$ in coordinate $V$.
The logarithmic specialization and amalgamation maps of
Stix, Theorem 3.4.6 and equation (6.2.5)
\cite{BC_Stix2002}, give the genuine based prime-to-five
comparison, by the same root-fiber-functor construction as F.9.
The outer component subgroup is $\langle d,c\rangle$.
A path based at $2$ inside this smooth outer component has
trivial inertia action in each prime-to-five quotient by based
good reduction.  The gluing along $XV=5$ therefore makes tame
inertia act by
\[
 a\longmapsto d^r a d^{-r},\qquad
 b\longmapsto d^r b d^{-r},\qquad c\longmapsto c,
 \quad r\in\prod_{q\ne5}\mathbb Z_q(1).
\]
Wild inertia acts trivially in this geometric prime-to-five
comparison.  These formulas follow also directly from the node
Kummer covers: a turn in the base annulus conjugates the bubble
loops by the common edge loop and fixes the based outer loops.
In the free basis $(x,y,z)=(b,d,c)$ the action is $\alpha_r$.

The image of $\Delta$ in tame inertia is
$4\prod_{q\ne5}\mathbb Z_q(1)$.  Indeed the tame part of the
five-adic cyclotomic character is the character of order four;
on wild inertia its image is $1+5\mathbb Z_5$, as is seen on
the extensions generated by fifth-power roots of unity.
A tame element in the kernel of the character of order four
can thus be corrected by a wild element to have cyclotomic
character one.  All other cyclotomic characters are already
trivial on inertia.  Projecting to the $\ell$ factor gives
$4\mathbb Z_\ell(1)$, including $\ell=2$.
It follows that the fixed group is the fixed group of
$\alpha_4$; continuity gives equality for the closure of its
integer powers.  Statement F.22 applies since $4\ne0$ in
$\mathbb Z_\ell$.  This proves (F.23) and its precise relation
to the full fixed group.
\end{proof}

\begin{statement}{F.24}{Proved}{The pro-five fixed groups for the two good-marked-reduction examples}
For $U=\mathbb P^1_{\mathbb Q_5}\setminus\{0,1,\infty\}$,
let $q_5:\pi\twoheadrightarrow\pi^{(5)}$ be the continuous
maximal geometric pro-five quotient.  With the genuine action
of the point $2$, respectively of an integral unit tangent at
$0$, we have
\[
 \bigl(\pi^{(5)}\bigr)^{s_2(\Delta)}=1,
 \qquad
 \bigl(\pi^{(5)}\bigr)^{s_{0,u}(\Delta)}=I_0^{(5)}.
                                                               \tag{F.24.1}
\]
Consequently $q_5(\pi^{s_2(\Delta)})=1$ and
$q_5(\pi^{s_{0,u}(\Delta)})=I_0^{(5)}$.
The latter equality of images follows because the full cusp
inertia is itself fixed.  It does not identify either full
profinite fixed group.
\end{statement}
\begin{proof}
Put $F=\mathbb Q(\zeta_5)$ and $S=\{v:v\mid10\}$, and let
$G$ be the maximal pro-five quotient of $G_{F,S}$.
We first prove that the image of inertia at five in $G$ is all
of $G$.  The discriminant of $F$ is $5^3$.  Its Minkowski
bound is $15\sqrt5/(2\pi^2)<2$, so its ideal class group is
trivial.  The prime two is inert, since $2$ has order four
modulo five, and $(\mathcal O_F/(2))^\times$ has order $15$.
The subgroup $\mu_5\subset\mathcal O_F^\times$ maps injectively
to this group and exhausts its five-primary part.
The ray-class exact sequence therefore gives
$\mathrm{Cl}_{(2)}(F)[5]=0$.
These uses of the ray-class correspondence are precisely
Milne, Chapter V, Theorem 3.6 and Example 3.9
\cite{BC_MilneCFT2020}.

A cyclic extension of degree five unramified outside $S$ and
unramified at five is tamely ramified at two and has conductor
dividing $(2)$.  The preceding computation excludes it.
If the inertia image were a proper closed subgroup of $G$,
its image in some finite five-group quotient would be proper
and would lie in a subgroup of index five.  Pulling back that
subgroup would give exactly such a cyclic extension.  This
proves the asserted surjectivity, without assuming that the
global group is free.

Let $F_S^{(5)}$ be the extension with group $G$.
Its intersection with $\mathbb Q(\mu_\infty)$ is
$F(\mu_{5^\infty})$.  To verify the equality, the intersection
is Galois and abelian over $\mathbb Q$, is unramified outside
$\{2,5\}$, and is pro-five over $F$.  The cyclotomic conductor
description over $\mathbb Q$ (Milne, Chapter V, Example 3.10)
shows that its part from a power of two has degree a power of
two; hence it is already in $F$.  The remaining part is the
five-power cyclotomic tower, which is contained in
$F_S^{(5)}$.  Thus the image of the global full-cyclotomic
kernel in $G$ is $\ker\chi_5$.
The local full-cyclotomic kernel has the same image: inertia
at five surjects onto $G$, and the cyclotomic character of an
inertia lift of an element of $\ker\chi_5$ is one at five and
at every other prime.  All these maps are restrictions of
Galois automorphisms for a fixed embedding
$\overline{\mathbb Q}\hookrightarrow\overline{\mathbb Q}_5$.

For the point $2$, the based pro-five action is unramified
outside $\{2,5\}$: the marked curve and its point extend over
$\mathbb Z[1/10]$, and based prime-to-residue-characteristic
specialization applies at every remaining prime.  Over $F$
its action on the Frattini quotient is trivial, since that
action is $\chi_5$ modulo five.  The kernel of
$\operatorname{Aut}(F_5(2))\to\operatorname{GL}_2(\mathbb F_5)$
is pro-five: the successive characteristic exponent-five
central quotients give finite filtrations with elementary
abelian five-group kernels.  Hence the based action factors
through $G$.  We have proved equality of the local and global
cyclotomic-kernel images for this actual based action.

We check that the global point section at $2$ is not cuspidal
even in the pro-five group.  Its two abelian Kummer coordinates
are $(\kappa(2),\kappa(-1))=(\kappa(2),0)$.
The maps $x,1-x:U\to\mathbb G_m$, evaluated on the leading
terms of the unit tangents, show that the cusp-one packet has
first coordinate zero and the cusp-infinity packet has equal
coordinates.  Since $v_{(2)}(2)=1$, neither contains this class.
To exclude the cusp-zero packet, consider over $F$ the
successive covers
\[
 y^5=1-x,\qquad
 z_j^5=\frac{y-\zeta_5^j}{-1-\zeta_5^j}\quad(j=1,2).       \tag{F.24.2}
\]
They are finite etale covers of the successive punctured
curves, with geometric degrees powers of five.  Their
composite is geometrically connected: the divisors of the two
displayed Kummer functions have independent classes modulo
five at $y=\zeta_5$ and $y=\zeta_5^2$.
The point lifts to $(y,z_1,z_2)=(-1,1,1)$.
Every cusp above zero first has $y=\zeta_5^i$.  Choose
$j\in\{1,2\}$ distinct from $i$.
At the prime $(1-\zeta_5)$ the numerator
$\zeta_5^i-\zeta_5^j$ has valuation one, while
$-1-\zeta_5^j$ is a unit.  This value is not a fifth power
in $F$, so no cusp above this point in (F.24.2) is $F$-rational.
An $F$-section contained in the lifted cusp decomposition group
would force such a rational cusp in every cover to which the
section lifts.  Thus the point section is noncuspidal.

The proof of B.4 applies word for word to this free pro-five
fundamental group over the number field $F$: its cover
comparison uses finite five-groups, its boundary sequence
uses $\mathbb Z_5(1)$, and Ribet's torsion finiteness applies
to the compactified cover Jacobians.  It gives a trivial global
cyclotomic-kernel centralizer for the just verified noncuspidal
point section, and cusp inertia for the unit tangent at zero.
Equality of the local and global images proves the first
formula in (F.24.1) and the second for tangent one.

Finally the pro-five completion of $\mathbb Z_5^\times$ is
free of rank one and is generated by the class of $2$:
$v_5(2^4-1)=1$, so $\log(2)$ generates $5\mathbb Z_5$.
An integral unit tangent $u$ therefore has
$\kappa_5(u)=a\kappa_5(2)$ locally for a unique $a\in\mathbb Z_5$.
The global cusp section obtained by twisting the unit tangent
by $a\kappa_5(2)$ is unramified outside $\{2,5\}$ and has the
same local pro-five section class.  It factors through the
same relative pro-five action considered above, and B.4 gives
its global centralizer $I_0^{(5)}$.  The local/global image
equality proves the claim for every integral unit tangent.
Conjugating a representative by cusp inertia leaves the
asserted inertia subgroup unchanged.  The stated maps from
the full fixed groups now follow directly by equivariance.
The exact finite certificates used here---the cyclotomic
discriminant, residue-field order, ray-class five-part,
Kummer-function norms, and the lifted point---are evaluated
in \texttt{code/bc\_local\_checks.py}, with its output included.
\end{proof}

\begin{statement}{F.25}{Proved}{The pro-five fixed group for the colliding-cusp example}
For $U=\mathbb P^1_{\mathbb Q_5}\setminus\{0,1,5,\infty\}$,
with the section at $2$ and
$\Delta=G_{\mathbb Q_5(\mu_\infty)}$, we have
\[
       \bigl(\pi^{(5)}\bigr)^{s_2(\Delta)}=1.                \tag{F.25.1}
\]
Consequently the full profinite fixed group maps to one
under the maximal pro-five quotient.  Together with
F.23 this evaluates every individual pro-prime fixed
quotient in all three examples, with the explicit
quotient maps described there.
\end{statement}
\begin{proof}
Filling the cusp $5$ is an actual based morphism
\[
 U\hookrightarrow U_0=\mathbb P^1\setminus\{0,1,\infty\}.
\]
It induces a continuous, surjective, equivariant map
$F_5(3)\to F_5(2)$ killing the normal closure $K$ of
the inertia at the new puncture.  By F.24 the target
has trivial fixed group for the section at $2$.
It suffices to prove $K^\Delta=1$.

We give the integral kernel comparison used for that
assertion.  Choose a geometric free basis which
expresses $F_5(3)=F_5(2)*_5\langle b\rangle$, where
$b$ is inertia at $5$ and the quotient kills $b$.
The profinite set $X$ of lifts of the point $5$ in
the universal pro-five cover of $U_0$, based at $2$,
is a torsor under $F_5(2)$.
After choosing one lift, identify it with $F_5(2)$.
The kernel $K$ is the free pro-five group on this
profinite set, with continuous generators
$x\mapsto xbx^{-1}$.
This free group is universal for continuous maps
from that profinite set to pro-five groups.
The assertion follows by two explicit inverse
homomorphisms: let the free group on $X$ be acted on
by $F_5(2)$ through left translation of $X$.
Its semidirect product with $F_5(2)$ receives the
free product by sending $b$ to the generator at
$1$ and fixing $F_5(2)$.
Conversely, $x\mapsto xbx^{-1}$ gives the homomorphism
back, equivariant for that left-translation action.
Both composites fix the indicated generators.
Thus the kernel and all its generator labels are
the actual ones for the filling morphism.

In particular
\[
 K^{\rm ab}\simeq\mathbb Z_5[[X]](1),                       \tag{F.25.2}
\]
where the completed module is the module of
$\mathbb Z_5$-valued finitely additive measures on
the clopen subsets of $X$.
Galois transports each inertia generator to a
conjugate of the generator at the transported
lifted cusp, raised to its cyclotomic power.
Thus (F.25.2) is equivariant, and on $\Delta$ the
twist is one.
For every $r\ge1$ there is also an injective
equivariant degree-$r$ Magnus homomorphism
\[
 \gamma_rK/\gamma_{r+1}K
          \hookrightarrow\mathbb Z_5[[X^r]].               \tag{F.25.3}
\]
To verify this integral comparison for the
infinitely generated kernel, use the finite
partitions $X\to X_n=F_5(2)/M_n$, with $M_n$
characteristic and open.  Every continuous map
from $X$ to a finite five-group factors through
one such partition, so the free group on $X$ is
the inverse limit of the free pro-five groups
on the finite sets $X_n$.
Its closed lower central groups and their
quotients are the corresponding inverse limits:
membership is detected in finite five-group
quotients, and compatible cosets lift by compactness.
For finite $X_n$, the Hall commutator basis
and its leading Magnus monomials give the
injective map into $\mathbb Z_5[X_n^r]$.
Taking their inverse limit gives (F.25.3).
Naturality of the degree-$r$ term makes this
map equivariant: conjugating a generator changes
only terms of degree at least two, and the action
on the degree-$r$ graded piece is induced from
its action on the degree-one module.

The action of $\Delta$ on $X$ has pro-five image.
Indeed its action on $F_5(2)$ is trivial on the
Frattini quotient, and that automorphism kernel
is pro-five, as proved in F.24.  After a choice
of torsor origin, the affine action on $X$ is
a homomorphism into the pro-five semidirect
product of $F_5(2)$ with this automorphism kernel.

There are no finite orbits of this action.
The regular function $1-x$ on $U_0$ sends its
path torsor from $2$ to $5$ to the division-root
torsor of
\[
                       \frac{1-5}{1-2}=4.
\]
On $\Delta$ its translation character is
$\kappa_5(4)=2\kappa_5(2)$.
The latter character has a unit value.
For a direct proof, the polynomial
$(T+2)^5-2$ is Eisenstein at five, so
$\mathbb Q_5(2^{1/5})$ has degree five.
It is not Galois: a Galois extension containing
one fifth root of $2$ contains all of them
and therefore $\zeta_5$, whose degree is four.
It cannot be a subfield of the abelian extension
$\mathbb Q_5(\mu_\infty)$.
Some element of $\Delta$ consequently moves that
fifth root, which says that $\kappa_5(2)$ has
nonzero reduction modulo five on $\Delta$.
Since it is a continuous homomorphism there,
its image is all of $\mathbb Z_5$.
Every orbit on $X$, and on $X^r$ by projection
to the first coordinate, therefore maps onto
an infinite translation orbit.

We finish with the required measure argument.
Let a pro-five group act continuously on a
profinite set $Z$ with no finite orbits.
An invariant integral $\mathbb Z_5$-measure
on $Z$ is zero.
For a clopen subset $C$, its stabilizer is open,
by uniform continuity on the compact set.
All orbits of this stabilizer on $C$ are
infinite.  For each $n$, choose a finite
invariant partition of $C$ on which every
orbit has size at least $5^n$.
Such a partition exists: each point has
arbitrarily many distinct translates separated
by a finite clopen partition; continuity gives
an invariant refinement, and a finite common
refinement works for all points by compactness.
The finite orbit sizes are powers of five.
Invariance makes the measure constant on the
atoms in each orbit, so its value on $C$ is
divisible by $5^n$.  Since this holds for every
$n$, that value is zero.  This proves the claim.

Apply it to $Z=X^r$ in (F.25.3).
Every graded fixed subgroup is zero.
A nontrivial element of $K$ has a nonzero image
in some lower-central graded quotient, because
it survives in a finite five-group and finite
five-groups are nilpotent.  A fixed element
would give a nonzero graded invariant, which
is impossible.  Thus $K^\Delta=1$, proving
(F.25.1) and the stated consequence for the
full profinite quotient map.
\end{proof}

\begin{statement}{F.26}{Proved}{A mixed-prime quotient and every one of its nilpotent fixed groups}
For the first example, with the section at $2$, consider the actual
finite etale cover
\[
 \begin{gathered}
 E^\circ\longrightarrow U,\qquad
 E:\ Y^2=X^3-2X,\quad t=X^2/2,\\
 D_E=\{O,(0,0),(\sqrt2,0),(-\sqrt2,0)\},
 \qquad P=(2,2).
 \end{gathered}
 \tag{F.26.1}
\]
Here $E^\circ=E\setminus D_E$, and $P$ maps to $2$.
The cover is cyclic of degree four over $\mathbb Q_5$:
in the coordinate $Z=Y/2$ its equation is
$Z^4=t(t-1)^2/2$. Choose $i\in\mathbb Q_5$ with
$i^2=-1$ and $i\equiv2\bmod5$; a deck generator is
$\delta(X,Y)=(-X,iY)$.
Let $H\subset\pi$ be the index-four subgroup determined by this
based cover, put $H_5=H^{(5)}$, and define
\[
 J=\pi/\ker(H\longrightarrow H_5),\qquad
 J_c=J/\gamma_{c+1}H_5\quad(c\ge1).
 \tag{F.26.2}
\]
These are continuous, surjective, $s_2(G_{\mathbb Q_5})$-equivariant
quotients of the full geometric profinite group. Their fixed groups
are
\[
 J^{s_2(\Delta)}=1,\qquad
 J_1^{s_2(\Delta)}=H_5^{\rm ab}\simeq\mathbb Z_5^5,
 \qquad
 J_c^{s_2(\Delta)}
   =\gamma_cH_5/\gamma_{c+1}H_5\quad(c\ge2).
 \tag{F.26.3}
\]
In particular the fixed group at class two is $\mathbb Z_5^{10}$,
and every nonzero element of it fails to have a fixed lift at
class three. The transition maps between consecutive fixed groups
in (F.26.3) are zero. The image of the full fixed group in $J$
is therefore trivial; in particular its image under the cyclic
quotient whose peripheral vector is $(1,2)\bmod4$ is trivial.
\end{statement}
\begin{proof}
We give the arithmetic and Lie comparisons explicitly.
The four points above $2$ are $P,\delta P,\delta^2P,\delta^3P$,
all rational over $\mathbb Q_5$. Thus the based arithmetic cover
is the one with section at $P$, and its subgroup $H$ is stable
under the genuine action in the statement. The kernel of the
maximal pro-five quotient of $H$ is characteristic in $H$,
normal in $\pi$, and stable under this arithmetic action.
This proves all the quotient assertions in (F.26.2).
The geometric group $H_5$ is free pro-five of rank
$2\cdot1+4-1=5$.

Write $V_E=H_5^{\rm ab}$ and
$W_E=\mathbb Z_5(1)[D_E]/\mathbb Z_5(1)\sum_d d$.
The functions
\[
             X,\quad X-\sqrt2,\quad X+\sqrt2
\]
have respective divisors $2d-2O$.
Their induced homomorphism on $V_E$ restricts to multiplication
by two on the boundary basis
$e_0,e_+,e_-$. Since two is a unit in $\mathbb Z_5$, it gives
an integral Galois-equivariant splitting
\[
                 V_E=W_E\oplus T_5E.                       \tag{F.26.4}
\]
The homomorphisms are equivariant over $\mathbb Q_5(\sqrt2)$;
their permutation under its Galois involution proves the stated
equivariance over $\mathbb Q_5$. They are also equivariant for
$\delta$, which fixes $e_0$ and interchanges $e_+$ and $e_-$.
Thus $W_E=W_+\oplus W_-$, with ranks two and one,
$W_-=\mathbb Z_5(e_+-e_-)(1)$.
Its arithmetic characters are respectively $\chi_5$ and
$\chi_5\varepsilon$, where $\varepsilon$ is the unramified
quadratic character. Indeed $\sqrt2$ generates the unramified
quadratic extension.

The two eigenspaces of $[i]$ split $T_5E$ integrally, since
$2i$ is a unit. One has an unramified character $\psi$;
the other has character $\chi_5\psi^{-1}$ by the Weil pairing.
Here is a direct check of the unramified assertion. The
endomorphism $[2+i]$ has degree five and unit differential
$2+i\bmod5$ on the good integral model. Its reduction is
separable, so its kernel and the kernels of its iterates are
finite etale over $\mathbb Z_5$. Their inverse limit is one
of the two CM Tate lines. Counting points gives
$\#E(\mathbb F_5)=10$, so the Frobenius polynomial is
$T^2+4T+5$. The unit root is $-2-i$ for the chosen embedding
of $i$; thus $\psi(\operatorname{Frob}_5)=-2-i$.
This root is not a root of unity: either complex embedding
has absolute value $\sqrt5$.
In particular $\Delta$ acts trivially on all of (F.26.4).

Choose the unit tangent at $O$, with parameter $X/Y$,
as an auxiliary reference section on $E^\circ$.
The deck automorphism has derivative $i$ there.
The compatible fifth-power roots of $i$ are all $i$,
and are rational over $\mathbb Q_5$. The resulting tangent
comparison gives a genuine order-four automorphism of the
based pro-five group, commuting with its arithmetic action.
Let
\[
 j_\sigma:V_E\longrightarrow\bigwedge^2V_E
 \quad(\sigma\in\Delta)
 \tag{F.26.5}
\]
be the first Johnson homomorphism: for a lift $\widetilde v$
of $v$, it is the class of
$\sigma(\widetilde v)\widetilde v^{-1}$ modulo $\gamma_3H_5$.
It is independent of that lift because $\Delta$ acts trivially
on $V_E$ and on $\bigwedge^2V_E$.
It is additive in $\sigma$ and is equivariant under conjugation
by $G_{\mathbb Q_5}$. The latter assertion follows directly
by conjugating an automorphism which is the identity on the
abelianization; any central part of the conjugating
automorphism cancels. It also commutes with $[i]$.

Each $d-O$, for $d\in D_E$, is two-torsion. Its pro-five
Kummer path is trivial: multiplication by $5^n$ fixes this
two-torsion point, providing a rational compatible division
sequence. The formula for the action on peripheral inertia
therefore gives
$j_\sigma(W_E)\subset\bigwedge^2W_E$.
The CM eigenvalues on $T_5E$ are $i,-i$, whereas those on
$\bigwedge^2W_E$ and on $\bigwedge^2T_5E$ are $1,-1$.
Commutation with $[i]$ consequently gives
$j_\sigma(T_5E)\subset W_E\otimes T_5E$.
Together these facts put (F.26.5), after tensoring with
$\mathbb Q_5$, in a direct sum of one-dimensional arithmetic
characters drawn from
\[
             \chi_5,\quad
             \varepsilon\psi^2,\quad
             \chi_5^2\varepsilon\psi^{-2}.                 \tag{F.26.6}
\]
For example the off-diagonal map from one CM line to
$W_-$ times the other has character
$\chi_5\varepsilon$ times the ratio of those two lines.
The diagonal maps use $W_+$ and have character $\chi_5$.
This proves (F.26.6) without assuming a diagonal action on
the nonabelian group.

Choose a geometric path from the auxiliary tangent to $P$,
and write $b_\sigma\in V_E$ for its abelian cocycle.
The change from the auxiliary action to the action at $P$
changes (F.26.5) by
\[
                    j_\sigma^P=j_\sigma+\operatorname{ad}b_\sigma;
                                                               \tag{F.26.7}
\]
reversing the initial path changes the two displayed signs
together and has no effect on the argument.
The functions defining (F.26.4), with leading coefficient
one at the auxiliary tangent, evaluate at $P$ as
$2,2-\sqrt2,2+\sqrt2$. Hence the component of $b$ in $W_-$ is
\[
 \frac14\kappa_5\!\left(\frac{2-\sqrt2}{2+\sqrt2}\right)
       (e_+-e_-)
    =\frac12\kappa_5(\sqrt2-1)(e_+-e_-).                    \tag{F.26.8}
\]
This is a cocycle for the character
$\theta=\chi_5\varepsilon$.
Its restriction to $\Delta$ is nonzero.
To prove that assertion, put $L=\mathbb Q_5(\sqrt2)$.
The unit $\sqrt2-1$ is not a root of unity, as its real
algebraic embedding has absolute value different from one.
The logarithm on a sufficiently small open unit subgroup
shows that its rational local Kummer class is nonzero.
For an abelian profinite group $\Gamma$ and a nontrivial
one-dimensional $\mathbb Q_5$-character $\eta$,
$H^1_{\rm cts}(\Gamma,\mathbb Q_5(\eta))=0$:
the cocycle identity for commuting $g,h$, with
$\eta(g)\ne1$, expresses every cocycle as the coboundary of
$c(g)/(\eta(g)-1)$.
Apply this to the abelian quotient $G_L/\Delta$ and $\chi_5$.
Inflation--restriction proves that the Kummer class stays
nonzero on $\Delta$.

The $T_5E$ component of $b|_\Delta$ is nonzero as well.
Indeed $\#E(\mathbb F_3)=4$ and $\#E(\mathbb F_5)=10$.
Injectivity of prime-to-residue-characteristic torsion
specialization bounds $E(\mathbb Q)_{\rm tors}$ by a group
of order dividing two. Since $Y(P)=2$, $P$ is not torsion.
The formal logarithm on an open subgroup of $E(\mathbb Q_5)$
then shows that its rational five-adic Kummer class is
nonzero. The local filtration and specialization assertions
are those of Milne, Chapter II, Theorem 4.1, Corollary 4.2
and Aside 4.4 \cite{BC_MilneEC2021}; alternatively the kernel
of reduction is a pro-residue-characteristic formal group,
which proves the torsion injectivity directly.
Both characters $\psi$ and $\chi_5\psi^{-1}$
are nontrivial. Applying the preceding abelian-quotient
argument to these two lines proves that this class also
stays nonzero on $\Delta$.

We now separate these homomorphisms on $\Delta$ exactly.
The character $\theta$ in (F.26.8) occurs neither in
(F.26.6), nor in $W_+$, nor in $T_5E$.
The powers of $\chi_5$ distinguish it from the two
off-diagonal characters in (F.26.6) on inertia; it differs
from $\chi_5$ by the nontrivial $\varepsilon$.
It differs from the unramified line $\psi$, and equality
with $\chi_5\psi^{-1}$ would require $\psi=\varepsilon$,
contradicting the displayed Frobenius value.
The joint homomorphism
\[
 \Delta\longrightarrow
  \operatorname{Hom}(V_E,\bigwedge^2V_E)\otimes\mathbb Q_5
           \ \oplus\ V_E\otimes\mathbb Q_5,\qquad
                 \sigma\longmapsto(j_\sigma,b_\sigma)
\]
has compact additive image, a $\mathbb Z_5$-lattice in
its rational span. That span is stable under
$G_{\mathbb Q_5}$ and is a sum of the characters just listed.
The projection onto its nonzero $\theta$ component is a
polynomial in finitely many of these commuting Galois
operators: for each other character choose an operator
distinguishing it from $\theta$, and multiply the corresponding
linear projections. Clearing the finitely many five-adic
denominators puts a nonzero projected vector back in the
compact image lattice. There is consequently an actual
$\sigma_0\in\Delta$ such that
\[
               j_{\sigma_0}=0,\qquad b_{\sigma_0}=w,
                    \qquad 0\ne w\in W_-\otimes\mathbb Q_5.
                                                               \tag{F.26.9}
\]
The vector in (F.26.9) is integral because it is in the
actual image. This step constructs the required group
element from the additive image, rather than choosing
unrelated values for different cocycles.

For completeness, the closed lower-central graded algebra
of the finite-rank free pro-five group embeds by its leading
Magnus terms in the tensor algebra on $V_E$.
An automorphism acting trivially on $V_E$ acts to the next
degree by the derivation induced by its first Johnson map.
This follows by expanding a commutator and using the
Leibniz rule in the tensor algebra; the Hall basis proves
the assertion integrally in each degree.
By (F.26.7)--(F.26.9), that derivation for $\sigma_0$
is $\operatorname{ad}w$.
In the free Lie algebra over $\mathbb Q_5$, the centralizer
of a nonzero degree-one element $w$ consists of its linear
span. One direct proof extends $w$ to a vector-space basis:
a homogeneous tensor polynomial commuting with the letter
$w$ contains only that letter, by comparing first and last
letters of its words. Its only possible homogeneous form
is a multiple of $w^r$; for $r>1$ this is not a Lie
element, as its tensor coproduct has nonzero intermediate
binomial coefficients.

It follows that the leading degree of a fixed element below
the last central degree of a nilpotent quotient must be one,
with leading vector proportional to $w$.
Choose $\sigma_1\in\Delta$ with $(b_{\sigma_1})_T\ne0$.
For this vector $w$, the $W_E\otimes T_5E$ component of
$j_{\sigma_1}^P(w)$ is exactly
$[(b_{\sigma_1})_T,w]$, which is nonzero.
Indeed $j_{\sigma_1}(w)$ and the bracket with
$(b_{\sigma_1})_W$ lie in $\bigwedge^2W_E$.
Thus no such degree-one vector is fixed either.
Every fixed element of $H_5/\gamma_{c+1}H_5$, for $c\ge2$,
is therefore in its last graded subgroup.
Conversely that subgroup is fixed because the action on
every associated graded is induced from the trivial action
on $V_E$. This proves the asserted fixed groups inside
$H_5$. Their ranks in degrees two and three are respectively
$10$ and $40$.

Finally no element outside $H_5$ can be fixed in $J_c$
for any $c\ge1$. The image of the inertia at $O$ in $J$
has an order-four subgroup mapping isomorphically onto the
deck group; its generator $\delta_0$ is fixed by the
auxiliary cuspidal action on $\Delta$.
Its conjugation on $V_E$ is the deck action $[i]$.
For an element with deck value $\delta_0^k$, the fixed
equation modulo $[H_5,H_5]$ for the action at $P$ requires
\[
                    (1-\delta^k)b_\sigma=0
                            \quad\hbox{for all }\sigma\in\Delta.
\]
For $k=1,2,3$, the operator $1-\delta^k$ is invertible
on $T_5E\otimes\mathbb Q_5$, and the proper component
of $b$ is nonzero. This excludes all three deck values.
Together with the preceding paragraph this proves the
formulas for $J_c$.
The intersection of the lower central subgroups of $H_5$
is one, since each nontrivial element survives in a finite
five-group, and every such finite group is nilpotent.
Taking the inverse limit proves $J^\Delta=1$.

The peripheral vector follows directly from the valuations
of $t(t-1)^2/2$ at $0$ and $1$.
All finite curve counts, toric identities, character
comparisons and Hall-dimension checks used here are
verified in \texttt{code/bc\_mixed\_cm\_checks.py}.
\end{proof}

\begin{statement}{F.27}{Proved}{An explicit finite obstruction to prime-to-five fixed lifts}
For the point $2$ on the tripod over $\mathbb Q_5$, the
image of the full fixed group in the geometric abelian
quotient $(\mathbb Z/4)^2$ is zero.
In particular, although every element of
$\pi^{(5')}$ is fixed, any such element with nonzero
abelian image modulo four has no fixed lift in $\pi$.
This obstruction is already detected by two actual finite
quotients, each isomorphic to
\[
                   (\mathbb F_5)^5\rtimes C_4
                   \quad\text{of order }12500.             \tag{F.27.1}
\]
The cyclic generator acts on the normal vector group by
fixing $e_0$, interchanging $e_+,e_-$, and multiplying the
two CM coordinates by $2$ and $3$.
\end{statement}
\begin{proof}
The first finite quotient is
$\overline J=J/(H_5^5[H_5,H_5])$, with the actual maps
of F.26. We verify at this precise finite level that
no nontrivial deck value is fixed.
Put $\beta=[2-i]$. Its kernel is the formal CM subgroup
of $E[5]$, and it is isomorphic over $\mathbb Q_5$ to
$\mu_5$: the other CM line has unramified character
$\psi\equiv1\bmod5$, and the Weil pairing identifies
the two lines as dual with cyclotomic twist.
The point $P=(2,2)$ is not in $\beta E(\mathbb Q_5)$.
Here is a finite certificate retaining the isogeny map.
On the integral affine curve it is given by
\[
 \begin{split}
 X(\beta Q)&=
  \frac{X(3+4i)(X^2-2-4i)^2}{(5X^2-2+4i)^2},\\
 Y(\beta Q)&=
  \frac{Y(2+11i)(X^2-2-4i)
          (X^4+(4+16i)X^2+4)}{(5X^2-2+4i)^3}.
 \end{split}                                                \tag{F.27.2}
\]
These identities follow by applying the addition law
to $2Q-[i]Q$, and are checked symbolically by the program.
The denominators are units at five. Modulo five the map
is $-\operatorname{Frob}_5$, so any preimage of $P$ must
reduce to $-P=(2,3)$. Modulo twenty-five, with $i=7$,
all five lifts of that residue point map under (F.27.2)
to $(7,2)$, rather than $(2,2)$. The finite enumeration
is included in the program and its output.

Consequently the isogeny Kummer class
$\delta_\beta(P)\in H^1(\mathbb Q_5,\mu_5)$ is nonzero.
It remains nonzero on $\Delta$.
Indeed for any $u\in\mathbb Q_5^\times$ which is not a
fifth power, adjoining a fifth root gives degree five:
after removing a fifth power and a fourth root of unity,
either its valuation is nonzero modulo five, and the
Newton polygon proves this degree, or it is
$(1+5)^b$ with $b\in\mathbb Z_5^\times$, in which case
$(1+T)^5-u$ is Eisenstein at five.
That degree-five extension is not Galois, since a
Galois extension containing one root would contain
$\zeta_5$, of degree four.
It cannot lie in the abelian extension
$\mathbb Q_5(\mu_\infty)$.
This proves the asserted restriction nonvanishing.
The map $[2+i]:E[5]\to E[\beta]$ sends the usual
multiplication-by-five Kummer class of $P$ to
$\delta_\beta(P)$. Thus the proper component of
$b|_\Delta$ in F.26 has nonzero reduction modulo five.

For $k=1,2,3$, $1-\delta^k$ is invertible on the
proper two-dimensional vector group over $\mathbb F_5$:
its eigenvalues are $1-2^k$ and $1-3^k$.
The last fixed equation in F.26 therefore excludes
nontrivial deck values already in $\overline J$.
Conversely the normal vector group is fixed because
$\Delta$ acts trivially on $V_E$.
Thus $\overline J^\Delta=(\mathbb F_5)^5$.
The order-four subgroup supplied by cusp inertia
splits the extension and has exactly the displayed
matrix, proving (F.27.1).

Write $(a,b)$ for the two abelian coordinates modulo
four in the original tripod. The cyclic cover equation
forces $a+2b=0$ for every fixed element.
The Mobius automorphism $m(t)=t/(t-1)$ is defined
over $\mathbb Q_5$ and fixes the base point $2$.
It therefore gives a genuine based automorphism of
$\pi$, commuting with the arithmetic action.
It fixes the cusp $0$ and interchanges $1,\infty$.
Its action on abelian inertia is
$e_0\mapsto e_0$, $e_1\mapsto-e_0-e_1$.
Pulling the first finite quotient back along this
actual automorphism gives the second quotient in
(F.27.1), with cyclic character $(1,1)\bmod4$.
Thus a fixed element also satisfies $a+b=0$.
Subtracting the two equations gives $b=0$ and then
$a=0$.

The abelian quotient modulo four factors through
$\pi^{(5')}$, by its universal property. Good marked
reduction with the integral point $2$ makes
$\pi^{(5')}$ fixed, as in F.23.
Therefore a prime-to-five element with nonzero image
in $(\mathbb Z/4)^2$ is an explicit fixed quotient
element having no fixed lift. All coordinate identities,
the five residue lifts and the invertible deck matrices
are checked in \texttt{code/bc\_mixed\_cm\_checks.py}.
\end{proof}

\begin{statement}{F.28}{Proved}{All abelian coordinates of the three full local fixed groups}
In the three examples of F.21, the images of the full
fixed groups in the full geometric abelianizations are,
respectively,
\[
 0,\qquad \widehat{\mathbb Z}\,e_0,\qquad 0.                 \tag{F.28.1}
\]
In particular the two point-section fixed groups lie in
the closed commutator subgroup of their respective full
geometric fundamental groups.

For every integer $m\ge2$ prime to five, the first
example has two explicit finite quotient tests excluding
all nonzero abelian coordinates modulo $m$.
They have the form
\[
 J_{m,n_m}=(\mathbb Z/5^{n_m})^{m+1}\rtimes C_m,
                                                               \tag{F.28.2}
\]
and their fixed groups are exactly their displayed
normal vector groups.  The integer $n_m$ is given by
the finite local valuation calculation (F.28.5) below.
Analogous finite quotients detect the second and third
images in (F.28.1).  The cases $m=2,3,4,6$ in the
first example all have $n_m=1$, giving orders
$250,1875,12500,468750$.
\end{statement}
\begin{proof}
Write $\Delta=G_{\mathbb Q_5(\mu_\infty)}$.
First consider the tripod with base point $2$ and
the actual cyclic cover
\[
 Y_m=\mathbb P^1_z\setminus(\{0,\infty\}\cup\mu_m),
       \qquad t=z^m.                                      \tag{F.28.3}
\]
Choose $\alpha^m=2$.  Its fiber consists of
$\alpha\zeta_m^j$.  All these points and all punctures
are defined over a finite unramified extension $L$
of $\mathbb Q_5$.  Indeed a residue extension of
degree $d$ with $4m\mid5^d-1$ contains every root
of $X^m-2$ and every $m$th root of unity; their
simple roots lift by Hensel's lemma.
Thus $L\subset\mathbb Q_5(\mu_\infty)$ and $\Delta$
fixes these fibers.

The kernel of $\pi\to C_m$ is the actual full
geometric fundamental group of $Y_m$.
Pass from it to its maximal pro-five abelian
quotient $V_m$, a free $\mathbb Z_5$-module of
rank $m+1$.  The kernel of this passage is
characteristic in the open subgroup and normal
in $\pi$.  Thus it gives a genuine quotient
$J_m$ of $\pi$ fitting into
$1\to V_m\to J_m\to C_m\to1$.
Its cyclic action fixes the boundary vector at
zero and cyclically permutes the $m$ vectors
at $\mu_m$.
The image of inertia at the ramified cusp zero
is procyclic with a unique prime-to-five factor
of order $m$.  This supplies an actual splitting,
so $J_m\simeq V_m\rtimes C_m$.
Since the punctures are $L$-rational and
$\chi_5|_\Delta=1$, the action of $\Delta$ on
$V_m$ is trivial.

Suppose that a fixed element has deck value
$j\ne0$.  Lift its based loop through (F.28.3).
It gives a $\Delta$-fixed point of the abelian
path torsor from $\alpha$ to
$\alpha\zeta_m^j$.  This description also
holds after passing to $V_m$: path lifting
commutes with the Galois action and with the
abelian quotient.  The function $z-1$, a unit
on $Y_m$, maps this torsor to the Kummer
torsor of
\[
              u_j=\frac{\zeta_m^j\alpha-1}{\alpha-1}.
                                                               \tag{F.28.4}
\]
It would therefore make
$\kappa_L(u_j)|_\Delta=0$.

This Kummer class is nonzero even with rational
five-adic coefficients.  Choose the complex
embedding with $\alpha=2^{1/m}>1$.
Then
$|\zeta_m^j\alpha-1|^2
=(\alpha-1)^2+2\alpha(1-\cos(2\pi j/m))$
is strictly larger than $(\alpha-1)^2$.
Consequently $u_j$ is not a root of unity.
For a finite local field, the five-adic
logarithm on principal units has precisely
torsion kernel; its valuation component
and logarithm identify the rational Kummer
completion of $L^\times$.
Thus its Kummer class is nonzero.
Restriction to $\Delta$ is injective on
$H^1(G_L,\mathbb Q_5(1))$.
To see this directly, the quotient
$\Gamma=G_L/\Delta$ is abelian, and its
cyclotomic character is nontrivial.
For a continuous cocycle $c$ on $\Gamma$,
commutativity gives
\[
  (\chi(g)-1)c(h)=(\chi(h)-1)c(g).
\]
Choose $g$ with $\chi(g)\ne1$ and divide
in $\mathbb Q_5$.  This expresses $c$ as
a coboundary.  Inflation--restriction,
obtained from the usual cochain maps,
proves the asserted injectivity.
This contradicts the fixed path and
excludes every nonzero deck value.

There is an explicit finite precision for
this exclusion.  Choose $L$ unramified
of degree $d$ as above and put $q=5^d$.
All $u_j$ are units: neither
$\alpha$ nor any $\zeta_m^j\alpha$ reduces
to one, since its $m$th power reduces to
$2$.  Set
\[
 e_j=v_L(u_j^{q-1}-1)-1,\qquad
             n_m=1+\max_{1\le j<m}e_j,
       \qquad v_L(5)=1.                                   \tag{F.28.5}
\]
These are finite nonnegative integers.
The decomposition
$\mathcal O_L^\times=
 \mu_{q-1}\times(1+5\mathcal O_L)$
and the logarithm isomorphism on the
second factor prove that $e_j$ is
exactly the largest $e$ for which
$u_j\in L^{\times5^e}$.
Write $u_j=v_j^{5^{e_j}}$ with $v_j$
not a fifth power.  Its nonzero
Kummer class modulo five remains
nonzero on $\Delta$.
After removing a Teichmuller factor,
$v_j\in1+5\mathcal O_L$ has
$v_L(v_j-1)=1$, so
$(1+T)^5-v_j$ is Eisenstein.
The degree-five extension generated by
a root is not Galois: a normal field
containing one root would contain
$\mu_5$, of degree four over the
unramified field $L$.
It therefore cannot be contained in
the abelian extension
$\mathbb Q_5(\mu_\infty)/L$.
Hence $\kappa_L(v_j)|_\Delta$ is
nonzero modulo five, and
$\kappa_L(u_j)|_\Delta$ is nonzero
modulo $5^{n_m}$.
This proves the finite exclusion in
$J_m/5^{n_m}V_m$.
Its normal vector group is fixed,
proving the exact fixed group in
(F.28.2).  Formula (F.28.5) involves
only algebraic valuations and gives
a terminating finite calculation.

Let $(a,b)$ be full abelian coordinates
at $0,1$.  The covering (F.28.3) forces
$a=0\bmod m$ for every fixed element.
The based automorphism
$t\mapsto t/(t-1)$ fixes $2$ and
commutes with the arithmetic action.
Its abelian map is
$(a,b)\mapsto(a-b,-b)$.
Pulling back (F.28.3) along this
actual automorphism forces
$a-b=0\bmod m$.  Therefore both
coordinates vanish modulo every
$m$ prime to five.
F.24 makes their five-primary
coordinates zero as well.
This proves the first image in
(F.28.1).

For the unit tangent at zero, use
instead the actual cover
$t=1-z^m$.  The induced tangents
above zero lie at the points
$\zeta_m^j$, with
$dz/dt=-1/(m\zeta_m^{j(m-1)})$.
Choose the starting tangent at
$z=1$, so its parameter $z-1$
has value $-1/m$.
A fixed element with deck value
$j\ne0$ would give a fixed
abelian path to the induced
tangent at $\zeta_m^j$.
The same function $z-1$ sends
these two fibers to Kummer
parameters $-1/m$ and
$\zeta_m^j-1$, respectively.
Their ratio is
$m(1-\zeta_m^j)$.
Its complex absolute value is
at least
$2m\sin(\pi/m)\ge4$, so it
is not a root of unity.
It is a five-adic unit because
$m$ is prime to five and
$\zeta_m^j$ has nontrivial
prime-to-five order.
The preceding restriction and
finite-precision arguments apply
verbatim, with these ratios.
This excludes its $e_1$ coordinate
modulo every $m$ prime to five.
F.24 gives the same exclusion
for its five-primary part.
Conversely the full cusp-zero
inertia is fixed by $\Delta$,
because the full cyclotomic
character is one there.
Its abelian image is the
whole primitive line
$\widehat{\mathbb Z}e_0$.
This proves the second equality.

For the four-puncture curve,
filling the cusp $5$ gives the
tripod point $2$ by an actual
based morphism.  Thus a fixed
element has zero $e_0,e_1$
coordinates.  The other based
map
\[
 t\longmapsto (5-t)/4:
 \mathbb P^1\setminus\{0,1,5,\infty\}
       \longrightarrow
 \mathbb P^1\setminus\{0,1,\infty\}
\]
fills the original cusp zero
and sends the point $2$ to
$3/4$, the cusp $5$ to zero,
and the cusp $1$ to one.
At the tripod point $3/4$,
the covers $t=z^m$ and
$t=1-z^m$ have starting roots
$(3/4)^{1/m}$ and
$(1/4)^{1/m}$.
Both are positive real numbers
different from one.
The strict absolute-value
inequality used for (F.28.4)
holds equally for a positive
root less than one.
Their reductions have $m$th
powers $2$ and $4$ modulo five,
so the same unramified and
unit arguments apply.
These two covers force both
prime-to-five abelian coordinates
of a fixed loop at $3/4$ to be
zero.  Pullback gives the
remaining $e_5$ exclusion.
F.25 excludes every five-primary
coordinate on the four-puncture
curve.  This proves its zero
abelian image.

The code
\texttt{bc\_cyclic\_fixed.py}
works in exact unramified
rings modulo powers of five.
It computes (F.28.5) for
$m=2,3,4,6$, obtaining
$e_j=0$ for every $j$ in
each case.  It verifies the
finite orders stated above.
\end{proof}

\begin{statement}{F.29}{Proved}{A genus-zero mixed quotient with no fixed element}
In the cyclic fourth-root cover $t=z^4$ of the
tripod with base point $2$, let $H_5$ be the
maximal geometric pro-five group of
$\mathbb P^1_z\setminus\{0,\infty,\pm1,\pm i\}$,
and let $J$ be the resulting quotient of the
full base fundamental group.  Thus
$1\to H_5\to J\to C_4\to1$ and $H_5$ is free
pro-five of rank five.
For the genuine point-section action,
\[
 J^\Delta=1,\qquad
 (J/\gamma_{c+1}H_5)^\Delta
       =\gamma_cH_5/\gamma_{c+1}H_5
                         \quad(c\ge1).                    \tag{F.29.1}
\]
Every transition between consecutive displayed
fixed groups is zero.
\end{statement}
\begin{proof}
Take $i\in\mathbb Q_5$, $i\equiv2\bmod5$,
and $\alpha^4=2$.
The extension $L=\mathbb Q_5(\alpha)$
is unramified of degree four: $X^4-2$
is irreducible and separable over
$\mathbb F_5$.
Its Frobenius sends $\alpha$ to $i\alpha$.
The auxiliary section is the unit tangent
at $z=0$.  The order-four factor of its
inertia supplies a based deck element
$\delta$, fixed by the auxiliary Galois
action.  Equivalently, the change of
tangent under $z\mapsto iz$ has the
compatible rational fifth-power roots
$i^{1/5^n}=i$, so its based comparison
commutes with $G_{\mathbb Q_5}$.

Write
\[
 V=H_5^{\rm ab}
   =\mathbb Z_5 e_0\oplus
                    \bigoplus_{j=0}^3\mathbb Z_5 f_j
\]
with a Tate twist on every summand;
the last four basis vectors are the
inertias at $i^j$.
All punctures are $\mathbb Q_5$-rational,
so the auxiliary arithmetic action on
$V$ is the scalar $\chi_5$.
For $\sigma\in\Delta$ its first
Johnson map is
$j_\sigma:V\to\bigwedge^2V$.
Composition in the first two
nilpotent degrees proves that
$j$ is additive on $\Delta$
and that
\[
 j_{g\sigma g^{-1}}=\chi_5(g)j_\sigma.                      \tag{F.29.2}
\]

Let $b$ be the cocycle comparing the
point $2$ with the auxiliary tangent
in $J$.  Its deck component records
$g\alpha/\alpha$.
On $\Delta$ this component is zero,
so $b_\sigma$ has an abelian value
in $V$, again an additive map of
$\Delta$.  We retain the precise
Frobenius transport.
Choose $\varphi$ lifting arithmetic
Frobenius with $\chi_5(\varphi)=1$;
this is possible because the fifth
cyclotomic extension is totally
ramified, whereas $L$ is unramified.
On $G_L$, the functional dual to
$f_j$ sends $b$ to
$\kappa_L(\alpha-i^j)$:
the unit $z-i^j$ has tangent value
$-i^j$, whose pro-five Kummer
class is zero.
For $\sigma\in\Delta$,
\[
 \begin{split}
 b_j(\varphi\sigma\varphi^{-1})
 &=\kappa_{\varphi^{-1}(\alpha-i^j)}(\sigma)\\
 &=\kappa_{\alpha-i^{j+1}}(\sigma)
   =b_{j+1}(\sigma).
 \end{split}                                              \tag{F.29.3}
\]
The discarded factor is $-i$, with
its compatible rational fifth-power
roots.  In the right-deck convention
this is the actual permutation
$\delta f_j=f_{j-1}$, while
$\delta e_0=e_0$.
The same conclusion follows directly
from the nonabelian cocycle identity:
modulo $[H_5,H_5]$ the auxiliary
$\Delta$-action is trivial, so
$b_{g\sigma g^{-1}}
 =b_g\rho(g)(b_\sigma)b_g^{-1}$.
Its linear action is
$\chi_5(g)\delta^{\,j(g)}$,
where $g\alpha=i^{j(g)}\alpha$.
This proves the required covariance
with the actual point-base deck
translation.

We now exhibit two nonzero eigenspace
components of $b|_\Delta$.
In the unramified integer ring
$\mathbb Z_5[\alpha]$,
the following are the coefficients
of $\log(\alpha-i^j)/5$ modulo five,
in the basis $1,\alpha,\alpha^2,\alpha^3$:
\[
 \begin{array}{c|rrrr}
 j&1&\alpha&\alpha^2&\alpha^3\\ \hline
 0&0&0&2&1\\
 1&0&0&3&2\\
 2&0&0&2&4\\
 3&0&0&3&3
 \end{array}.                                             \tag{F.29.4}
\]
Here the logarithm kills the Teichmuller
factor.  This is an exact finite
calculation: with $i=7\bmod25$,
compute
$((\alpha-i^j)^{624}-1)/(5\cdot624)$
in
$(\mathbb Z/25)[\alpha]/(\alpha^4-2)$
and reduce modulo five.
The program performs these operations
and verifies every row.

For the permutation in (F.29.3),
the coefficient at $f_0$ of the
$\lambda$-eigenprojection on the
four puncture coordinates is
$\frac14\sum_j\lambda^{-j}b_j$.
For $\lambda=i$ and $\lambda=-1$,
(F.29.4) gives respectively
$\alpha^3$ and $2\alpha^2$ modulo
five after dividing the logarithm
by five.  Thus both rational
Kummer classes are nonzero.
The restriction injection proved
in F.28 makes their restrictions
to $\Delta$ nonzero as well.

The joint additive image of
$\Delta$ under $\sigma\mapsto
(j_\sigma,b_\sigma)$ is a compact
$\mathbb Z_5$-submodule.
Conjugation by $\varphi$ preserves
it, acts trivially on $j$ by
(F.29.2), and by the order-four
permutation $\delta$ on $b$.
The integral operators
\[
        E_\lambda=\frac14
             \sum_{a=0}^3\lambda^{-a}\varphi^a
                      \quad(\lambda=i,-1)
                                                               \tag{F.29.5}
\]
therefore preserve this actual image.
They annihilate its $j$-coordinate
and retain the indicated nonzero
$b$-component.  Compactness makes
the image closed under every
$\mathbb Z_5$ scalar in (F.29.5).
Consequently there are actual
$\sigma_1,\sigma_2\in\Delta$
with
\[
 j_{\sigma_1}=j_{\sigma_2}=0,\qquad
 b_{\sigma_1}=w_i\ne0,\quad
 b_{\sigma_2}=w_{-1}\ne0,                                  \tag{F.29.6}
\]
where the two vectors lie in
distinct deck eigenspaces.
They are linearly independent.

Changing from the auxiliary tangent
to the point section changes the
first Johnson map to
$j_\sigma+\operatorname{ad}b_\sigma$.
Thus (F.29.6) supplies two independent
pure-inner first Johnson directions.
The leading Magnus term of a fixed
element below the final degree of
a nilpotent quotient must commute
in the free Lie algebra with both
$w_i$ and $w_{-1}$.
As proved in F.26, the centralizer
of a nonzero degree-one vector in
a free Lie algebra over $\mathbb Q_5$
is its one-dimensional span.
The two spans intersect in zero.
The integral Hall graded modules
embed in their rationalizations,
so this excludes every such
nonzero leading term.
Conversely the final central
graded subgroup is fixed because
$\Delta$ acts trivially on $V$
and hence on every free Lie
graded quotient.

F.28 with $m=4,n_m=1$ excludes
all nonzero deck values already
in the finite abelian-normal
quotient.  It excludes them
in every quotient in (F.29.1)
as well.  This proves the
displayed finite-depth groups.
Their images at the preceding
depth are zero.  Finally
$\bigcap_c\gamma_cH_5=1$:
every nontrivial element of a
free pro-five group survives
in a finite five-group, and
finite five-groups are nilpotent.
Taking the inverse limit proves
$J^\Delta=1$.
\end{proof}

\begin{statement}{F.30}{Proved}{A torsion test in every actual finite cover}
Put $K_\infty=\mathbb Q_5(\mu_\infty)$ and
$\Delta=G_{K_\infty}$.
For every finite extension $L/\mathbb Q_5$ put
$\Delta_L=G_{LK_\infty}$ and
$\Gamma_L=\operatorname{Gal}(LK_\infty/L)$.
For every semiabelian variety $A/L$,
restriction is injective:
\[
 H^1(G_L,V_5A)\longrightarrow
 H^1(\Delta_L,V_5A),\qquad
 V_5A=T_5A\otimes_{\mathbb Z_5}\mathbb Q_5.
                                                        \tag{F.30.1}
\]
Consequently, for every smooth curve $Y/L$, its smooth
compactification $C$, reduced boundary $D=C\setminus Y$,
and $P,Q\in Y(L)$, a $\Delta_L$-fixed geometric
\'etale path from $P$ to $Q$ forces
\[
                 [Q-P]\in J_D(L)_{\mathrm{tors}}.
                                                        \tag{F.30.2}
\]
The same conclusion holds for a fixed path in the
maximal abelian pro-five path torsor.
When $L\subset K_\infty$ one has $\Delta_L=\Delta$.
Here $J_D$ is the generalized Jacobian with reduced
modulus $D$; when $D$ is empty it is the ordinary Jacobian.
For nonempty $D$, condition (F.30.2) has the exact
divisor formulation
\[
 \exists n\geq1,\ f\in\overline L(C)^\times:
 \quad \operatorname{div}(f)=n(Q-P),\qquad f|_D=1.
                                                        \tag{F.30.3}
\]

In particular, let $U/\mathbb Q_5$ have a rational
interior base point $x$, let
$\pi=\pi_1(U_{\overline{\mathbb Q}_5},x)$, and use
the point section to define the action of $\Delta$.
For every $\Delta$-stable open subgroup $H\subset\pi$,
the corresponding based cover over $K_\infty$ is an
actual finite \'etale cover.  If $h\in\pi^\Delta$,
its lift beginning at the chosen point $P$ ends at
a $K_\infty$-point $Q$ above $x$, and this pair satisfies
(F.30.2)--(F.30.3) after descent to a finite
$L\subset K_\infty$.
\end{statement}

\begin{proof}
We first prove the representation-theoretic assertion
used in (F.30.1).
Let $\Gamma$ be an abelian profinite group and $M$ a
finite-dimensional continuous $\mathbb Q_5$-representation
with $M^\Gamma=0$.
Let $R\subset\operatorname{End}_{\mathbb Q_5}(M)$ be the
unital algebra generated by $\Gamma$.
It is a finite-dimensional commutative algebra acting
faithfully on $M$.  The ideal
$I=(\gamma-1:\gamma\in\Gamma)$ is all of $R$.
Indeed, if it lay in a maximal ideal $\mathfrak m$,
the corresponding local Artinian factor $R_{\mathfrak m}$
would act on a nonzero summand $M_{\mathfrak m}$.
Its last nonzero radical power
$\mathfrak m^aM_{\mathfrak m}$ would contain a nonzero
vector killed by $\mathfrak m$, hence by every $\gamma-1$.
This contradicts $M^\Gamma=0$.
Thus there are finitely many $A_i\in R$ and
$\gamma_i\in\Gamma$ with
\[
                  \sum_i A_i(\gamma_i-1)=1.
\]
For a continuous one-cocycle $c:\Gamma\to M$,
commutativity gives
$(\gamma_i-1)c(\gamma)=(\gamma-1)c(\gamma_i)$.
Consequently
\[
 c(\gamma)=(\gamma-1)m,\qquad
                 m=\sum_iA_ic(\gamma_i).
                                                        \tag{F.30.4}
\]
This is an actual continuous coboundary formula;
in particular $H^1(\Gamma,M)=0$.

We record explicitly the local point-group input.
For a semiabelian variety $A$ over the finite local
field $L$, a sufficiently small formal neighborhood
of the identity is an open subgroup
$U_A\simeq\mathbb Z_5^d$.
The torus part has a finitely generated valuation
quotient, and the abelian-variety part is compact.
The local analytic subgroup assertion for the latter
is Mattuck's Theorem~7, printed p.~114
\cite{BC_Mattuck1955}.
These statements give a finitely generated quotient
$A(L)/U_A$.
For completeness, the image in the abelian-variety
part is open, since the differential of the smooth
quotient morphism is surjective; it is therefore
compact. Finitely many local analytic sections over
this compact image, together with generators for the
torus valuation lattice, give the asserted finite
generation modulo $U_A$.
Passing to the preimage of a torsion-free subgroup
of finite index in $A(L)/U_A$ and lifting its free
basis gives a finite-index subgroup of the form
\[
                       \mathbb Z^r\oplus\mathbb Z_5^d.
                                                        \tag{F.30.5}
\]
In particular $A(L)_{\mathrm{tors}}$ is finite and
$(V_5A)^{G_L}=0$.
It also follows directly from the finite Kummer
sequences and (F.30.5) that the map
\[
                   A(L)\longrightarrow H^1(G_L,V_5A)
                                                        \tag{F.30.6}
\]
has exactly the torsion subgroup as its kernel.
One way to see the last assertion without a completion
identification is to start with a zero rational Kummer
class. A five-power multiple of the point is then
divisible by every $5^n$ in $A(L)$.
Multiplication of this point and all its division
witnesses by the finite index of (F.30.5) places
them in that subgroup.
The intersection of the subgroups $5^n$ times
(F.30.5) is zero.

Now put $V=V_5A$ and $M=V^{\Delta_L}$.
The quotient $\Gamma=\Gamma_L=G_L/\Delta_L$ is
abelian: restriction identifies it with the closed
subgroup $\operatorname{Gal}(K_\infty/(L\cap K_\infty))$
of the abelian cyclotomic Galois group.
and
$M^\Gamma=V^{G_L}=0$.
Let a cocycle on $G_L$ represent an element in the
kernel of (F.30.1).  Subtracting an actual coboundary
makes its restriction to $\Delta_L$ zero.
The cocycle identity then shows both that its values
lie in $M$ and that it is constant on cosets of $\Delta_L$.
It descends continuously to $\Gamma$.
Formula (F.30.4) makes it a coboundary and proves
(F.30.1).

The path-to-Kummer map in the second assertion is the
one induced by the actual Abel morphism
\[
                  \iota_P:Y\longrightarrow J_D,\qquad
                    Q'\longmapsto[Q'-P].
                                                        \tag{F.30.7}
\]
More explicitly, pull back the finite \'etale isogenies
$[5^n]:J_D\to J_D$ along $\iota_P$.
A path from $P$ to $Q$ carries the compatible zero
points above $\iota_P(P)=0$ to compatible division
points of $[Q-P]$.
Its displacement under $\sigma$ is the cocycle
$\sigma Q_n-Q_n$ in $J_D[5^n]$.
Thus a $\Delta_L$-fixed path makes the restriction of
the rational Kummer class of $[Q-P]$ zero.
The construction factors through the abelian pro-five
path torsor because these deck groups are abelian
five-groups.
Apply (F.30.1) and then (F.30.6) to obtain
(F.30.2).
The quotient of degree-zero divisors disjoint from $D$
by divisors of functions equal to one on $D$ is the
actual generalized-Jacobian group.  This gives
(F.30.3); see also \cite[\S2.1, equation (2.2),
and \S2.3, Lemma~2.5]{BC_Creutz2019}.

Finally the arithmetic subgroup for the indicated
cover is
$H\rtimes s_x(\Delta)\subset\pi\rtimes s_x(\Delta)$.
It has finite index, and its chosen fibre point is
$\Delta$-fixed.  The usual lifting map for finite
\'etale covers identifies paths in $Y$ with base
paths having the prescribed endpoints.
The lift of $h$ is fixed because both $h$ and its
starting point are fixed; its endpoint is therefore
a $K_\infty$-point.
The cover, the two points, and their morphism to $U$
are given by finitely many coefficients in
$K_\infty$, so they descend together to some finite
$L\subset K_\infty$.
The previous construction then proves the stated
test in this actual cover.
Further extension to any finite $L'/L$ preserves the
fixed-path conclusion, since
$\Delta_{L'}=G_{L'K_\infty}\subset\Delta_L$.
Thus this passage also permits arbitrary finite
extensions used to change the reduction of the cover.
\end{proof}

\begin{statement}{F.31}{Proved}{The full abelian fixed images for all genus-zero curves}
Let $L/\mathbb Q_5$ be finite,
$K_\infty=\mathbb Q_5(\mu_\infty)$, and
$\Delta_L=G_{LK_\infty}$.
Let $Y/L$ be a smooth genus-zero hyperbolic curve
and $P\in Y(L)$.
For the action defined by this actual point section,
\[
 \pi_1(Y_{\overline L},P)^{\Delta_L}
       \longrightarrow \pi_1(Y_{\overline L},P)^{\mathrm{ab}}
       \quad\hbox{has image zero}.                       \tag{F.31.1}
\]
There is no restriction on the reduction of $Y$,
on the position of $P$, or on the prime divisors of
the finite abelian quotient being tested.
For an $L$-rational nonzero tangential base point
$\vec v$ at an $L$-rational cusp $c$, the corresponding
full fixed group has abelian image exactly
$\widehat{\mathbb Z}(1)e_c$.

The same conclusion applies at every actual
genus-zero finite \'etale neighborhood of $P$
whose chosen fibre point is fixed by $\Delta_L$,
after a finite extension of $L$.
Thus a full fixed element has zero abelian image in
every such neighborhood in which it lifts to a loop.
\end{statement}

\begin{proof}
Write $C$ for the compactification and $D=C\setminus Y$.
The rational point $P$ identifies $C$ with a projective
line.  Extend $L$ finitely to split $D$; a fixed element
remains fixed by the smaller group $\Delta_{L'}$.
We retain the letter $L$ for this extension during the
construction.

Choose a cusp $d_\infty\in D$, and an actual
$L$-isomorphism
$z:C\longrightarrow\mathbb P^1_L$
sending $d_\infty$ to infinity.
Put $p=z(P)$.
For each $d\in D\setminus\{d_\infty\}$ set
\[
       u_d=\frac{z-z(d)}{p-z(d)}:Y\longrightarrow\mathbb G_m.
                                                        \tag{F.31.2}
\]
These maps give the usual full abelian coordinates:
the inertia at $d$ maps to $1$ in
$\widehat{\mathbb Z}(1)$, the inertia at
$d_\infty$ maps to $-1$, and every other cusp inertia
maps to zero.  The sole relation among the cusp
generators is their sum, so the product of these
induced maps is an isomorphism
$\pi_1(Y_{\overline L},P)^{\mathrm{ab}}
 \simeq\widehat{\mathbb Z}(1)^{|D|-1}$.

Fix $n\geq3$ and form the actual connected cyclic cover
\[
  Y_{d,n}\longrightarrow Y,\qquad
              w^n=u_d,\qquad
              z=z(d)+(p-z(d))w^n.                        \tag{F.31.3}
\]
Its compactification is the projective $w$-line.
The removed points are $0,\infty$ and the $n$-th
roots of
$(z(e)-z(d))/(p-z(d))$ for
$e\in D\setminus\{d,d_\infty\}$.
This lists precisely all inverse images of $D$,
and proves that (F.31.3) is finite \'etale on
the indicated open curves.
Its fibre above $P$ is $\mu_n$.
Every one of these points is fixed by $\Delta_L$.

Let $h$ be a full fixed element whose image under
$u_d$ modulo $n$ is $j\ne0$.
Lift its actual path to (F.31.3), beginning at $w=1$.
The endpoint is $w=\zeta_n^j$, for a compatible
choice of generator $\zeta_n$.
The path is fixed by $\Delta_L$.
Choose another cusp
$e\in D\setminus\{d,d_\infty\}$, which exists by
hyperbolicity.  Put
\[
       r=\frac{z(e)-z(d)}{p-z(d)}\ne0,1,\qquad \beta^n=r.
\]
After a finite extension $L'/L$ containing $\beta$
and $\mu_n$, each
\[
                 v_a=w-\beta\zeta_n^a
                       \quad(0\leq a<n)                 \tag{F.31.4}
\]
is an actual unit on $Y_{d,n}$.
It is nonzero at both endpoints, since $r\ne1$.
The induced map of paths to $\mathbb G_m$ sends the
fixed path to the Kummer path of
\[
             R_a=
       \frac{\zeta_n^j-\beta\zeta_n^a}
            {1-\beta\zeta_n^a}\in L'^\times.             \tag{F.31.5}
\]
Formula (F.30.1), now with
$\Delta_{L'}=G_{L'K_\infty}\subset\Delta_L$,
therefore forces every $R_a$ to be a root of unity.

Here is an exact contradiction.
Embed the finitely generated characteristic-zero
field containing these finitely many coefficients
and roots into $\mathbb C$.
Set
\[
 D_a=|\zeta_n^j-\beta\zeta_n^a|^2
                   -|1-\beta\zeta_n^a|^2.
\]
The root-of-unity sums
$\sum_a\zeta_n^{2a}=\sum_a\zeta_n^{-2a}=0$,
valid for $n\geq3$, give
\[
 \sum_{a=0}^{n-1}D_a^2
       =2n\,|\zeta_n^j-1|^2\,|\beta|^2>0.                \tag{F.31.6}
\]
Indeed
$D_a=-(\zeta_n^j-1)\overline\beta\,\zeta_n^{-a}
      -(\zeta_n^{-j}-1)\beta\,\zeta_n^a$,
and squaring proves the displayed identity.
Thus some $|R_a|\ne1$, so that $R_a$ cannot be a
root of unity.  This excludes every $j\ne0$.

We have proved that every coordinate in (F.31.2)
vanishes modulo every $n\geq3$.
Every positive modulus divides one of these $n$,
so all full abelian coordinates vanish.
The integers $n$ used in (F.31.3) may be divisible
by five: the maps are \'etale in characteristic zero,
and the proof uses the actual subgroup
$\Delta_{L'}$, which fixes all roots of unity.
This proves (F.31.1).

For the tangential assertion, choose
$d_\infty\ne c$ and use the same actual coordinate
$z$, putting $p=z(c)$.
For each cusp $d\ne c,d_\infty$ use
(F.31.3) with this value of $p$ and $n\geq5$.
The cover is unramified at $c$.
If the original tangent has velocity $\dot z$,
its lift at $w=1$ has velocity
$\dot z/(n(p-z(d)))$; the deck translate at
$w=\zeta_n^j$ has velocity
$\zeta_n^j\dot z/(n(p-z(d)))$.
These are the actual derivatives of (F.31.3),
so both lifted tangents are fixed by $\Delta_{L'}$.
The other points $\zeta_n^a$ above $c$ are punctures.
For $a\ne0,j$, the unit $w-\zeta_n^a$ has
nonzero value at both tangential endpoints.
Its induced Kummer ratio is therefore
$(\zeta_n^j-\zeta_n^a)/(1-\zeta_n^a)$;
no tangent velocity enters this nonvanishing value.
Put $A=|\zeta_n^j-1|^2$.
Formula (F.31.6) with $\beta=1$, after subtracting
the terms at $a=0,j$, gives
\[
 \sum_{\substack{0\leq a<n\\a\ne0,j}}
 \left(|\zeta_n^j-\zeta_n^a|^2
                  -|1-\zeta_n^a|^2\right)^2
                      =2A(n-A)>0,                       \tag{F.31.7}
\]
because $0<A\leq4<n$.
The identical Kummer argument excludes all $j\ne0$.
This kills every abelian coordinate except the
one corresponding to $c$.
Conversely the full inertia
$I_c\simeq\widehat{\mathbb Z}(1)$ at this actual
tangent is fixed by $\Delta_L$:
its action on compatible roots of the local
parameter is the cyclotomic character, which is
one on $\Delta_L$.
Its image is precisely the primitive cusp line.

The obstruction can also be made finite in each
instance. For an $R_a$ selected by (F.31.6), its
restricted integral Kummer homomorphism
$\Delta_{L'}\to\mathbb Z_5(1)$ is nonzero, since its
rational class is nonzero. It remains nonzero modulo
some $5^s$. Pulling back the cover
$T^{5^s}=v_a$ along (F.31.4) supplies a finite
\'etale cover on which the indicated endpoints
cannot be joined by a fixed path.
Taking its normal closure, and then the compositum
of the finitely many coefficient-field conjugates,
gives an actual finite geometric quotient stable
under the original arithmetic action.  The conjugate
family is finite because all the coefficients,
points, roots, and covering equations used here
descend to the finite extension $L'/L$.

Finally apply the same construction to the given
finite neighborhood, using its chosen fixed fibre
point as $P$. Descent of the cover and point to a
finite field and subsequent restriction of
$\Delta_L$ are exactly the morphisms described at
the end of F.30.
The exact root-of-unity identities (F.31.6)--(F.31.7)
and the cyclic-cover inverse maps are checked by
\texttt{bc\_genus\_zero.py}.
\end{proof}


\begin{statement}{F.32}{Proved}{The full local fixed group at every point and tangent}
Let $L/\mathbb Q_5$ be finite, let
$K_\infty=\mathbb Q_5(\mu_\infty)$, and put
$\Delta_L=G_{LK_\infty}\subset G_L$.
Let $X/L$ be a smooth geometrically connected
hyperbolic curve, and let $\Pi$ be its full
geometric profinite fundamental group.
For an $L$-rational point $b\in X(L)$ and its actual
point section $s_b$,
\[
 C_\Pi(s_b(\Delta_L))=1.                                  \tag{F.32.1}
\]
For a nonzero $L$-rational tangent at an
$L$-rational cusp $c$ and its actual tangential
section $s_{\vec c}$,
\[
 C_\Pi(s_{\vec c}(\Delta_L))=I_c,                          \tag{F.32.2}
\]
where $I_c$ is the actual geometric inertia
subgroup selected by that tangent.
In the tripod example of Task F.4,
$I_0\simeq\widehat{\mathbb Z}(1)$, with this
isomorphism induced by the compatible local
covers $w^n=t$.
\end{statement}
\begin{proof}
We use the resolution of nonsingularities theorem
of Mochizuki--Tsujimura, in the author's revision
of 30 March 2024.  The exact inputs are
Theorem 2.17 and Proposition 2.4(vii):
with $\Sigma$ the set of all primes, a section
of a closed subgroup $H\subset G_L$ fixes at most
one primitive valuation in their universal
geometric covering, provided $H\cap I_L$ has a
continuous quotient $\mathbb Z_\ell$ for some
$\ell\ne5$.  The correspondence between valuations
and center-chains, including its inverse, is
Proposition 2.3(v)--(vi)
\cite[Propositions 2.3(v)--(vi), 2.4(vii);
Theorem 2.17]{BC_MochizukiTsujimura2024}.
We verify the subgroup hypothesis and construct
the valuation and its stabilizer explicitly.

First every finite unramified extension of $L$
is obtained by adjoining suitable roots of unity
of order prime to five.  Indeed, such roots
give all finite extensions of the finite residue
field, and their lifts are the unique roots
supplied by Hensel's lemma.
Consequently $L^{\mathrm{ur}}\subset LK_\infty$
and
\[
 \Delta_L\subset I_L,\qquad
 I_L/\Delta_L
   \ \xrightarrow[\ \chi_5\ ]{\ \sim\ }\
 \chi_5(I_L)\subset\mathbb Z_5^\times
                 \simeq C_4\times\mathbb Z_5.             \tag{F.32.3}
\]
For the kernel assertion, inertia already fixes
all prime-to-five roots of unity, so its full
cyclotomic kernel is exactly its five-adic
cyclotomic kernel.  Thus (F.32.3) is induced by
the actual cyclotomic homomorphism.

Choose compatible three-power roots of a
uniformizer of $L$.  The resulting tame
homomorphism $t_3:I_L\to\mathbb Z_3(1)$ is onto:
the successive extensions over the maximal
unramified field have degrees $3^n$, as follows
either from their valuations or from the
Eisenstein polynomials for those roots.
Let $B=t_3(\Delta_L)$, a closed subgroup.
The quotient $\mathbb Z_3(1)/B$ is a continuous
quotient of $I_L/\Delta_L$.
Every finite quotient of the latter has order
divisible only by two and five, whereas every
finite quotient of $\mathbb Z_3(1)/B$ is a
three-group.  Therefore $B=\mathbb Z_3(1)$.
This verifies the uniqueness hypothesis with
$H=\Delta_L$ and $\ell=3$.
It uses the at-most-one part of Proposition
2.4(vii), whose hypothesis imposes no nontrivial
cyclotomic image on $H$.

Here is the promised geometric comparison.
Write $M$ for the union of the function fields
of the connected finite geometric covers of
$X_{\overline L}$ inside its chosen universal
cover.  Then
\[
 \Pi=\operatorname{Gal}(M/\overline L(X)).
\]
More precisely, transport the chosen based fibre
functor to a geometric generic point by an
\'etale path.  The transported arithmetic
fundamental group is
$\operatorname{Gal}(M/L(X))$; its projection to
$G_L$ is restriction to the constant field,
with kernel the displayed group $\Pi$.
Transport the section and the basepoint lifts
by the same path.  Thus every action below is
an action in this actual extension.
The based cover construction gives a compatible
lift of $b$ to every finite cover.
At a tangent, the chosen compatible Puiseux
branches give instead a compatible cusp in
the normalizations of the smooth compactification.
In either case write $\widetilde b$ for this
compatible collection of centers.
Let
\[
 B_{\widetilde b}
   =\bigcup_Y\mathcal O_{\overline Y,b_Y}
           \subset M,\qquad
 \epsilon_{\widetilde b}:B_{\widetilde b}
                  \longrightarrow\overline L            \tag{F.32.4}
\]
be the directed union of the local discrete
valuation rings at the centers and its residue
map.  In (F.32.4), $\overline Y$ is the smooth
compactification, so the formula applies to
cusps as well as to interior points.
The union is a valuation ring: every element
of $M^\times$ lies in a finite subcover, where
it or its inverse belongs to the indicated
local discrete valuation ring.
Its residue field is $\overline L$, because
all the centers are geometric and the constant
field already contains $\overline L$.

Compose this horizontal valuation with the
five-adic valuation of its residue field.
The resulting valuation ring is
\[
 R_{\widetilde b}
   =\epsilon_{\widetilde b}^{-1}
                       (\mathcal O_{\overline L})
                      \subset B_{\widetilde b}.           \tag{F.32.5}
\]
Its intersection with $L$ is $\mathcal O_L$.
It is primitive in the precise sense of
Mochizuki--Tsujimura Definition 2.2(i):
its maximal ideal is the only prime ideal
containing five.
To prove this directly, put
$J=\ker\epsilon_{\widetilde b}$.
Since $5^{-n}J=J$ for every $n\ge0$, any prime
ideal of $R_{\widetilde b}$ containing five
contains $J$.
The quotient $R_{\widetilde b}/J$ is the
rank-one valuation ring $\mathcal O_{\overline L}$,
whose only prime ideal containing five is its
maximal ideal.  Taking its inverse image proves
the assertion.

The morphism from this construction to the
source theorem is the following actual map:
send $R_{\widetilde b}$ to its center in every
compactified semistable model of every finite
subcover, and then to the vertex or edge
containing that center.
These entries form the center-chain of
Definition 2.2(v).
Proposition 2.3(v) identifies this map with a
bijection from the valuation rings in question
to the center-chain space; Proposition 2.3(vi)
gives its inverse as the directed union of the
local rings at the centers.
Both maps commute with the arithmetic
fundamental-group action, since formation of a
center and pullback of a local ring commute
with the corresponding field automorphisms.
In particular equality of the primitive
center-chains here is equality of the rings
$R_{\widetilde b}$.

The actual section $s_b(G_L)$ preserves the
compatible basepoint lifts, while the actual
section $s_{\vec c}(G_L)$ preserves the
compatible cusps selected by the tangent.
Their action on the residue field in (F.32.4)
is the ordinary Galois action on $\overline L$,
which preserves $\mathcal O_{\overline L}$.
Thus $R_{\widetilde b}$, and hence its primitive
center-chain, is fixed by the restricted
section $s(\Delta_L)$.

Now take $h\in C_\Pi(s(\Delta_L))$.
For every $\sigma\in\Delta_L$,
\[
 s(\sigma)\bigl(hR_{\widetilde b}\bigr)
       =h\,s(\sigma)R_{\widetilde b}
       =hR_{\widetilde b}.
\]
Thus $hR_{\widetilde b}$ is a second fixed
primitive valuation.  The uniqueness theorem
gives $hR_{\widetilde b}=R_{\widetilde b}$.

We next recover the exact geometric stabilizer,
without identifying different kinds of
valuation by notation.
There is an equality of subrings of $M$
\[
 R_{\widetilde b}[1/5]=B_{\widetilde b}.                    \tag{F.32.6}
\]
Indeed five is a unit in the horizontal
valuation ring.  Conversely, for any
$f\in B_{\widetilde b}$, a sufficiently large
power $5^nf$ has integral residue and belongs
to $R_{\widetilde b}$.
It follows that an element of $\Pi$ preserving
$R_{\widetilde b}$ preserves
$B_{\widetilde b}$, hence the center $b_Y$
on every finite normalization $\overline Y$.
Conversely, an element preserving all these
centers preserves $B_{\widetilde b}$ and its
residue map, since it fixes $\overline L$
pointwise; it therefore preserves
$R_{\widetilde b}$.
We have proved the exact stabilizer equality
\[
 \operatorname{Stab}_\Pi(R_{\widetilde b})
       =\operatorname{Stab}_\Pi(\widetilde b).              \tag{F.32.7}
\]

At an interior point, every finite geometric
Galois cover is \'{e}tale there.
Its deck group acts freely on the fibre:
the residue extension is trivial over
$\overline L$ and the ramification index is one.
Thus an element of (F.32.7) has identity image
in every finite geometric Galois quotient,
and is itself the identity.  This proves
(F.32.1).

At a cusp, (F.32.7) is the inverse limit of the
inertia groups at the chosen compatible
normalization points, which is exactly the
image $I_c$ of the punctured local disc.
For completeness, each completed local
extension has the form
$\overline L((t^{1/e}))/\overline L((t))$:
its residue field is algebraically closed,
its residue characteristic for this
horizontal valuation is zero, and an
$e$-th root of the unit multiplying a local
uniformizer exists by Hensel's lemma.
Its automorphisms multiply $t^{1/e}$ by
$\mu_e$.
These compatible homomorphisms identify the
stabilizer with the actual local inertia image,
including all transition maps.
For the tripod at zero the global finite
covers $w^n=t$ realize every such cyclic
quotient, so the local map
$\widehat{\mathbb Z}(1)\to\Pi$ is injective
and has image $I_0$.

Finally the tangential section acts on these
compatible local roots by the cyclotomic
character.  Since $\Delta_L$ fixes all roots
of unity, it acts trivially on $I_c$.
The inclusion $I_c\subset
C_\Pi(s_{\vec c}(\Delta_L))$, together with
(F.32.7) and uniqueness, proves (F.32.2).
\end{proof}
\begin{statement}{F.21}{Proved}{The three full local fixed groups and all fixed lifts}
For the three examples in Task F.4, with
$\Delta=G_{\mathbb Q_5(\mu_\infty)}$ and the actual
section action, the full profinite fixed groups are
\[
 \begin{array}{c|c|c}
 U&\text{section}&\pi^{s(\Delta)}\\ \hline
 \mathbb P^1\setminus\{0,1,\infty\}&2&1\\
 \mathbb P^1\setminus\{0,1,\infty\}
        &\text{integral unit tangent at }0
        &I_0\simeq\widehat{\mathbb Z}(1)\\
 \mathbb P^1\setminus\{0,1,5,\infty\}&2&1 .
 \end{array}                                               \tag{F.21.1}
\]
For $q:\pi\to\pi^{(5')}$, the images of these
full fixed groups are, respectively,
\[
            1,\qquad I_0^{(5')},\qquad 1.                  \tag{F.21.2}
\]
Consequently the finite coboundary tests of F.20
all hold exactly for $\bar h=1$ in the two
point-based examples, and exactly for
$\bar h\in I_0^{(5')}$ in the tangential example.
The fixed lift is unique in each point-based
example.  Each tangential lift fibre is a coset
of $\mathbb Z_5(1)\subset I_0$.
\end{statement}
\begin{proof}
Formula (F.21.1) follows from F.32 applied to the
two actual interior point sections and the
actual tangent at zero.
On the inertia selected by this tangent,
the quotient $q$ is the actual projection
\[
 \widehat{\mathbb Z}(1)
      =\prod_\ell\mathbb Z_\ell(1)
      \longrightarrow
        \prod_{\ell\ne5}\mathbb Z_\ell(1).
\]
The covers $w^n=t$ with $(n,5)=1$ verify
injectivity of the displayed target into
$\pi^{(5')}$ and identify it with
$I_0^{(5')}$.
Its kernel is the five-adic factor.
This proves (F.21.2) and the stated lift fibres.
Finally F.20 identifies the existence of each
such fixed lift with all the finite coboundary
tests, proving the exact evaluation of those
tests as well.

The earlier F.22--F.31 computations retain their
explicit quotient actions and finite obstruction
witnesses.  The universal-cover valuation
argument in F.32 supplies the detection beyond
those particular quotient families.
\end{proof}

